\documentclass[11pt]{article}
\usepackage[margin=1in]{geometry}
\usepackage{amsmath,amssymb,amsthm}
\usepackage[T1]{fontenc}
\usepackage{lmodern}
\usepackage[hidelinks]{hyperref}

\title{Question 1 of arXiv:2408.08935 is answered by arXiv:2602.01421}
\author{fa\_banach\_001 literature-status packet}
\date{2026-06-18}

\begin{document}
\maketitle

\paragraph{Status.}
This is a \emph{literature already answered} packet.  It records an exact
later answer to one open question from Andrea Garcia, \emph{Greedy algorithms:
a review and open problems}, arXiv:2408.08935.

\paragraph{Source question.}
In Section 3 of arXiv:2408.08935, after defining the Power-Relaxed Greedy
Algorithm with relaxation parameter \(m^{-\alpha}\), the source asks Question
1: whether it is possible to find \(\alpha_0>1\) such that
\[
  \|f-\mathcal T_m^r(f)\| \lesssim m^{-\alpha_0/2},\qquad m=1,2,\ldots,
\]
for every \(f\in A_1(\mathcal D)\).  This asks whether the exponent in the
DeVore--Temlyakov relaxed greedy estimate can be improved by taking
\(\alpha>1\) in the power-relaxed variant.

\paragraph{Supporting answer.}
Pablo M. Berna and Andrea Garcia, \emph{Convergence Analysis of Greedy
Algorithms with Adaptive Relaxation in Hilbert Spaces}, arXiv:2602.01421,
explicitly restates this source problem as Question 2.3, citing
``Greedy algorithms: a review and open problems (2025), Question 1''.  Their
Theorem 2.4 gives a negative answer: for every \(\alpha>1\), there is a
dictionary in \(\mathbb R^2\) and an element \(f\in A_1(\mathcal D)\) for which
the Power-Relaxed Greedy Algorithm fails to have the desired convergence.  The
proof uses the fact that
\[
  P_\alpha=\prod_{k=2}^{\infty}\left(1-\frac1{k^\alpha}\right)
\]
is strictly positive when \(\alpha>1\), yielding a positive lower bound on the
error.

\paragraph{Scope limitations.}
This identification answers only Question 1 of arXiv:2408.08935.  It does not
answer Question 2 on alternative step functions, nor Questions 3--5 concerning
semi-greedy bases, the isometric \(1\)-almost-greedy versus \(1\)-semi-greedy
problem, or semi-greedy renormings.

\paragraph{References.}
\begin{itemize}
\item A. Garcia, \emph{Greedy algorithms: a review and open problems},
  arXiv:2408.08935.
\item P. M. Berna and A. Garcia, \emph{Convergence Analysis of Greedy
  Algorithms with Adaptive Relaxation in Hilbert Spaces}, arXiv:2602.01421.
\end{itemize}

\end{document}

